\section{Introduction}
\label{sec:intro}

Weak magnetic fields, whose energy per atom is negligible compared with the
thermal energy $kT$, nevertheless change the plasticity of nonmagnetic
crystals; the broad class of such observations is known as the magnetoplastic
effect \cite{Alshits2003,Molotskii2000}. Aluminum alloys with Fe-bearing
inclusions occupy a special place in this class because, if such an inclusion
is ferromagnetic, they offer a seemingly straightforward mechanical route from
the field to plastic flow: the field magnetizes the inclusion;
magnetostriction---the change of shape that a magnetized body
undergoes---strains it; and the stress caused by the mismatch between the
strained inclusion and the surrounding matrix is transferred to the matrix.
Whether the inclusions in these alloys are in fact ferromagnetic is considered
in Section~\ref{sec:phase}. A series of experiments on a commercial
Al--Mg--Si alloy with iron-bearing microinclusions has documented the effect
in detail: a 30~min exposure to $B=0.4$--$0.7$~T, applied \emph{before}
mechanical testing, increases the subsequent room-temperature creep
strain---the slow deformation under a constant load---at 160~MPa by roughly a
quarter, increases the heat release during deformation several-fold, and
shifts the Taylor--Quinney coefficient, the fraction of the work of plastic
deformation that is released as heat, from 0.07 to 0.25
\cite{Skvortsov2019JAP,Skvortsov2019PSS,Friha2024JMMM,Nikolaev2025SciRep}.
Because the effect remains after the field is removed, it has been described
as ``magnetic memory'' in plasticity \cite{Skvortsov2019PSS}.

The interpretation advanced in these works rests on a single estimate. The
stress at the inclusion--matrix interface produced by magnetostriction of the
inclusion, obtained from the field strength through an empirical
proportionality coefficient, is put at about 147~MPa for $B=0.7$~T
\cite{Friha2024JMMM}. This exceeds the 120~MPa yield stress of the matrix,
and the conclusion drawn is that the matrix around every inclusion is
plastically deformed and its density of dislocations---the line defects of
the crystal lattice whose motion produces plastic flow---is raised. This chain
of reasoning has never been examined at the scale on which it is supposed to
act. Whether a stress of this magnitude actually appears in the matrix, over
what distance from the interface it persists, and whether it is capable of
nucleating dislocations or of setting existing ones in motion are questions
about {\aa}ngstr{\"o}ms and picoseconds---the native domain of classical
molecular dynamics (MD), in which the atoms move according to Newton's
equations under forces derived from an interatomic potential.

In the present work the magnetostrictive mechanism is subjected to a
quantitative atomistic test at the interface between aluminum and the
intermetallic compound Al$_{13}$Fe$_4$. A magnetic field has no direct
representation in classical MD, so its action is represented mechanically:
based on the earlier experimental work \cite{Friha2024JMMM}, the inclusion
is taken to elongate along the field direction by
$\varepsilon=1.94\times10^{-3}$ (0.194\%) at constant volume, the strain
corresponding to the 147~MPa interface stress estimated there. A direct
reproduction of the macroscopic effect in a simulation cell is not attempted;
the reason is set out in Section~\ref{sec:discussion}. Instead, each link of
the proposed chain is examined separately. The following quantities are
computed: the stress field that the elongated inclusion creates in the matrix
and its decay with distance from the interface; the applied shear stress at
which a dislocation already present in the matrix starts to move, and that
and whether new dislocations are nucleated at the interface up to the highest stress applied; and the resistance
that Mg and Si solute atoms dissolved in the aluminum oppose to dislocation
motion. The chain is then closed from the opposite end: a two-scale estimate
gives the interface stress that the measured 25\% creep enhancement
\emph{requires}, assuming thermally activated glide---dislocation motion in
which obstacles are overcome with the help of thermal fluctuations, so that
the glide rate depends exponentially on the stress---and the decay as
$r^{-3}$ of the elastic field outside a strained inclusion
\cite{Eshelby1957}.